\documentclass[12pt]{article}
\usepackage[utf8]{inputenc}
\usepackage[T1]{fontenc}
\usepackage[margin=1in]{geometry}
\usepackage{setspace}
\usepackage{xcolor}
\usepackage[colorlinks=true, urlcolor=blue!60!black, linkcolor=blue!60!black]{hyperref}
\setlength{\parindent}{0em}
\setlength{\parskip}{0.6em}
\renewcommand{\baselinestretch}{1.0}
\begin{document}
\begin{center}
{\LARGE \textit{Moore v. Harper}, 600 U.S.\ 1 (2023)\par}
\vspace{0.5em}
{\large Case Brief --- Two-Page Edition\par}
\end{center}

\section{Memory Jogger}
The Court rejected the ``independent state legislature theory'': state courts may review congressional maps under state constitutions, within the ordinary bounds of judicial review (\href{https://www.law.cornell.edu/supremecourt/text/21-1271}{Cornell LII}).

\section{Detailed Case Facts}
North Carolina's Republican legislature drew partisan congressional maps in 2021. Voters including Rebecca Harper sued in state court under the state constitution's Free Elections Clause. Speaker Timothy Moore argued that the Federal Elections Clause gives state legislatures exclusive, unchecked power over congressional maps.

\section{Procedural History}
The North Carolina Supreme Court struck the maps in \textit{Harper I}. After the U.S. Supreme Court granted certiorari, a newly Republican state court overruled \textit{Harper I} in \textit{Harper III} but left its injunction in place. The Supreme Court affirmed \textit{Harper I}.

\section{Judicial Votes}
Affirmed, 6--3. Roberts wrote for Sotomayor, Kagan, Kavanaugh, Barrett, and Jackson; Kavanaugh concurred; Thomas dissented, joined by Gorsuch and, on mootness, Alito.

\section{Holding}
The case was not moot. The Elections Clause does not free legislatures from state constitutional limits or state judicial review, but federal courts must ensure state courts do not exceed the ordinary bounds of judicial review. No test was adopted.

\section{Analysis of Opinions}
\textbf{Majority:} Judicial review was a Founding-era ``fundamental principle''; \textit{Smiley v. Holm} and \textit{Arizona State Legislature} treat election rules as lawmaking bound by state constitutions, and no procedure-substance line exists. \textbf{Kavanaugh:} adopt Rehnquist's \textit{Bush v. Gore} test, measured against pre-existing state law. \textbf{Thomas:} the case was ``indisputably moot,'' state constitutions may not impose substantive limits, and federal review will federalize state law.

\section{Examples: Future Applications}
\textbf{Same-Side:} a state court enforces express compactness rules against a map. \textbf{Opposite-Side:} days before an election, a court extends an unambiguous ballot deadline on ``policy fairness'' alone, inviting federal intervention. \textbf{Fence-Sitter:} a first-ever reading of a ``free elections'' clause to extend a deadline during a hurricane.

\section{Critique}
Originalists including Akhil and Vikram Amar and Steven Calabresi praised the rejection of the theory, while formalists echoed Thomas on mootness. Progressive scholars such as Leah Litman and Katherine Shaw warn that the undefined ``ordinary bounds'' limit is a ``time bomb'' making federal judges supervisors of state election law.

\section{Key Quotations}
\begin{quote}
``The Elections Clause does not insulate state legislatures from the ordinary exercise of state judicial review.''
\end{quote}
\noindent --- Chief Justice Roberts, Opinion of the Court.

\end{document}
